\section{长期价值的收益为什么要拆开看}
\subsection{先修数据可见性，再训练}
我要求保留真实数据身份，不能把签到代理称为真实订单。数据Agent重现教师25,550条local-clock代理后，Codex主线程又取两端UTC复算，发现12对时间矛盾、6对UTC倒序。Research独立确认后，保留教师主基线，另外剔除副本做稳健性，没有偷偷统一纽约时区或重新定义原数据。

另一项更直接的修正是司机热点。原Notebook会读全量签到，本次开发只读inner签到，完整重训后也只读前70\%日期。beta扫描在验证池完成，未来测试池只在方案冻结后用于确认。300回合是实际发生的模拟训练，不是300天真实运营；这一身份贯穿曲线、表格和报告。

\subsection{Double检验没有被简化为“修复成功”}
beta1下的收入下降伴随大量非正边权。Reinforcement提出并执行Double Q手算、小例和五训练seed比较，再由独立Evolution核验。Double提高了beta1相对收入，却降低应答，拒绝的可行边比例也未降低；所以没有把理论上的过估计修正直接写成业务效果全面改善。

过程中还发现了工程混杂：更新coin与行为探索共用RNG，Double多消耗随机数会改变后续行为流。执行器暂停受影响进程，保留旧Double身份，分离随机流后重跑；Q-only数组逐项bit等价的旧结果继续按原身份引用。这个闭环由真实问题触发，不是为了让过程报告看起来更曲折。

\subsection{从无成本调度到移除边权的回退}
主动调度首先完整模拟了相邻移动、距离、时间与忙时，司机没有瞬移。无门槛版本虽然增加gross收入，却产生约5000 km/模拟日的额外空驶。看到这一负结果后，Codex才提出扣除$\lambda d$的移动门槛，并单独比较0/1/3/5，说明价值/公里的教学量纲，不假定它是真实能源价格。

原Q-KM边权与价值调度组合并未自动形成最强策略。五训练seed下，移除边权LT、保留同一个learned V做移动，gross收入反而每次都更高，因此冻结KM即时派单+门槛3的调度作为gross候选。更省空驶的门槛5保留为替代；Double和NSGA的独立作用也没有被塞进这个组合中充数。

冻结后的首批5日，候选相对Q-KM的gross差为640.81，但区间跨0；随后预登记20个新采样seed，仅确认同一冻结方法，差为435.45。额外空驶848.08 km意味着按1单位/km的情景，收入减额外调度成本之差为$-412.62$，区间也为负。这条负结果改变了解释：可以说所测gross收入增加，不能说已证明真实净收益。

\subsection{融合产生的是权衡，不是自动升级}
在独立Q-KM和NSGA成立后，才用长期边权与接驾构造前沿，并把固定膝点/约束选点带回完整仿真。约束方案降低了接驾，但收入和响应也下降；膝点更进一步牺牲服务。因此融合完成了教师加分的实测要求，却没有被称为最终全面最优。

我在执行提示词中要求保留没有收益的合法实验，这一要求在这里得到落实。具体采用、移除和停止是Codex在批准规则内依据实测作出的执行判断；没有补写成我看到每一条结果后都亲自选择，也没有把另一个角色的反馈伪装成新的人工质疑。
